% Corrected French for printed p.146.
% Correction: opérateurs de bord / opérateur de bord total in footnote 6.
\SourcePage{146}

\LayerTarget{tohoku:example:bicomplex-spectral-sequences}
Soit par exemple $K$ un \textit{bicomplexe}\footnote{Contrairement à la
terminologie introduite dans [6], nous supposons que les deux opérateurs de
bord $d'$ et $d''$ d'un bicomplexe $K$ \textit{commutent}, et nous prendrons
donc comme opérateur de bord \textit{total} $d=d'+d_t''$, où
$d_t''x=(-1)^p d''x$ pour $x\in K^{p,q}$.} dans $C$, $K=(K^{p,q})$, et
supposons que pour tout entier $n$, il n'existe qu'un nombre fini de couples
$(p,q)$ tels que $p+q=n$ et que $K^{p,q}\ne0$. Alors on peut trouver deux
suites spectrales, aboutissant toutes deux à $H(K)=(H^n(K))$ (cohomologie du
complexe simple $(K^n)$ associé à $K$,
$K^n=\sum_{p+q=n}K^{p,q}$) et dont les premiers termes sont respectivement
\LayerTarget{tohoku:eq:2.4.1}
\[
\tag{2.4.1}
{}_{\mathrm I}E_2^{p,q}(K)=H_{\mathrm I}^pH_{\mathrm{II}}^q(K),
\qquad
{}_{\mathrm{II}}E_2^{p,q}(K)=H_{\mathrm{II}}^pH_{\mathrm I}^q(K)
\]
(en employant les notations de [6, Chap. XV, N\textsuperscript{o}6]). Ces
suites spectrales sont des foncteurs spectraux en $K$. Ce sont d'ailleurs des
suites spectrales \textit{cohomologiques} si les degrés de $K$ sont positifs.

\LayerTarget{tohoku:construction:right-derived-spectral-functors:p1}
Soit $F$ un foncteur covariant d'une catégorie abélienne $C$ dans une autre
$C'$, supposons que dans $C$ tout objet soit isomorphe à un sous-truc d'un
objet injectif. Soit $K=(K^n)$ un complexe dans $C$, limité à gauche (i. e.
$K^n=0$ pour $n$ assez petit). Alors les considérations de [6, Chap. XVII]
s'appliquent, et permettent de construire deux suites spectrales,
$({}_{\mathrm I}E_r^{p,q}(K),{}_{\mathrm I}F^n(K))$ et
$({}_{\mathrm{II}}E_r^{p,q}(K),{}_{\mathrm{II}}F^n(K))$, aboutissant au même
truc gradué $\mathfrak R F(K)=(\mathfrak R^nF(K))$ (pour deux filtrations
convenables de ce dernier), et dont les termes initiaux sont respectivement :
\LayerTarget{tohoku:eq:2.4.2}
\[
\tag{2.4.2}
{}_{\mathrm I}E_2^{p,q}(K)=H^p(R^qF(K)),
\qquad
{}_{\mathrm{II}}E_2^{p,q}(K)=R^pF(H^q(K)).
\]
(Comme d'habitude, un foncteur appliqué à un complexe $K$ désigne le complexe
obtenu en appliquant le foncteur à chacun des termes de $K$). Ces suites
spectrales sont des foncteurs par rapport à $K$, et la variance des termes
initiaux et finaux est celle qui est évidente sur leur forme explicite. On a
ainsi défini deux foncteurs spectraux sur la catégorie des complexes dans $C$
limités à gauche, on les appelle les \textit{foncteurs spectraux dérivés à
droite de $F$} ou encore les \textit{foncteurs spectraux d'hyperhomologie (à
droite) de $F$}. Les foncteurs $\mathfrak R^nF(K)$ sont appelés les
\textit{foncteurs d'hyperhomologie de $F$}.

\LayerTarget{tohoku:construction:hyperhomology-resolution:a:p1}
Rappelons le principe de la construction, en nous bornant pour fixer les idées
au cas où $K$ est à degrés positifs. Soit $L=(L^{p,q})$ un complexe double, à
degrés positifs, muni d'une augmentation $K\longrightarrow L$ (où $K$ est
considéré comme un bicomplexe où le second degré est nul) ; on suppose que pour
tout $p$, le complexe $L^{p,*}=(L^{p,q})_q$ est une \textit{résolution} de
$K^p$, et que pour tout $p,q$, on a $R^nF(L^{p,q})=0$ pour $n>0$. Voici deux
façons particulières, toutes deux uniques à une équivalence d'homotopie de
bicomplexes près, de construire un tel bicomplexe : a) On considère $K$ comme
un objet de la catégorie abélienne $\mathfrak K$ des complexes dans $C$ à
degrés positifs, et on prend pour $L$ une résolution \textit{injective} de
l'objet $K$. On constate facilement que les objets injectifs dans $\mathfrak K$
sont les complexes $A=(A^i)$ à degrés positifs tels que chaque $A^i$ soit
injectif, $H^i(A)=0$ pour $i>0$, et qui ``se décomposent'' (i. e. les
sous-trucs $Z(A^i)$ des cycles sont facteurs directs dans les $A^i$). De plus,
tout objet de $\mathfrak K$ se plonge dans

\TohokuContinuation{147}
